\documentclass[11pt,a4paper]{article}

% ── Packages ─────────────────────────────────────────────────────
\usepackage[margin=1in]{geometry}
\usepackage{amsmath,amssymb}
\usepackage{graphicx}
\usepackage{booktabs}
\usepackage{hyperref}
\usepackage{xcolor}
\usepackage{listings}
\usepackage{authblk}
\usepackage{abstract}
\usepackage{titlesec}
\usepackage{natbib}
\usepackage{microtype}
\usepackage{parskip}

% ── Hyperlink colours ─────────────────────────────────────────────
\hypersetup{
    colorlinks=true,
    linkcolor=blue!70!black,
    citecolor=blue!70!black,
    urlcolor=blue!70!black
}

% ── Title ─────────────────────────────────────────────────────────
\title{\textbf{Hallucination Fingerprints: Consistent Failure
Patterns in Large Language Models}}

\author[1]{Nikhil Upadhyay}
\affil[1]{Dublin Business School, Dublin, Ireland \\
\texttt{nikhil25000@gmail.com}}

\date{April 2026}

% ─────────────────────────────────────────────────────────────────
\begin{document}
\maketitle

% ── Abstract ──────────────────────────────────────────────────────
\begin{abstract}
Large language models hallucinate---producing factually incorrect
outputs with high confidence---in ways that are not fully understood.
We investigate whether hallucinations follow \emph{consistent,
detectable patterns} in the internal activations of transformer
models prior to the generation of an incorrect token.
Through experiments on a custom 806K-parameter transformer and
GPT-2 (124M parameters), tested across 20,000 factual prompts in
7 knowledge categories, we identify two named phenomena:
\textbf{Relation Dropout}, in which attention to the semantic
relation token collapses in the final transformer block of small
models; and \textbf{Last-Layer Suppression}, in which factual
knowledge emerges in blocks 10--11 of GPT-2 but is systematically
suppressed by block 12.
We further propose a three-type hallucination taxonomy
and release HallScan, an open-source detection tool, and
HallBench, a labeled benchmark of 20,000 annotated examples.
\end{abstract}

% ── 1. Introduction ───────────────────────────────────────────────
\section{Introduction}

Large language models (LLMs) have demonstrated remarkable
capabilities in knowledge retrieval, reasoning, and language
generation \citep{brown2020language}. However, they remain prone
to \emph{hallucination}: the confident generation of factually
incorrect content \citep{ji2023survey}. Despite extensive research
into hallucination detection at the output level, the internal
mechanisms that give rise to hallucinations remain poorly understood.

In this paper, we ask a more fundamental question: \emph{do
language models hallucinate in consistent, detectable patterns
visible in their internal activations before the incorrect token
is generated?} If such patterns exist, they could form the basis
of a new class of hallucination detectors that operate on model
internals rather than post-hoc output verification.

Our contributions are as follows:
\begin{itemize}
    \item We identify and name \textbf{Relation Dropout}: a
    phenomenon in which small transformer models lose attention
    to the semantic relation token in the final block immediately
    before a hallucination.

    \item We identify and name \textbf{Last-Layer Suppression}:
    a phenomenon in which GPT-2 correctly identifies factual
    answers in blocks 10--11 of its residual stream, but block
    12 systematically suppresses this signal via structural
    pattern interference.

    \item We propose a \textbf{three-type hallucination taxonomy}
    (Type 1: Ignorance/Dropout; Type 2a: Suppression; Type 2b:
    Knowledge Gap) grounded in internal activation analysis.

    \item We validate our findings at scale across \textbf{20,000
    factual prompts} in 7 knowledge categories using an RTX 4060
    GPU.

    \item We release \textbf{HallScan}, an open-source Python
    library for hallucination detection, and \textbf{HallBench},
    a labeled benchmark of 20,000 annotated examples.
\end{itemize}

% ── 2. Related Work ───────────────────────────────────────────────
\section{Related Work}

\paragraph{Hallucination in NLG.}
\citet{ji2023survey} provide a comprehensive survey of
hallucination in natural language generation, distinguishing
between intrinsic hallucinations (contradicting source material)
and extrinsic hallucinations (unverifiable claims).
\citet{lin2022truthfulqa} introduce TruthfulQA, a benchmark
measuring whether models mimic human falsehoods, finding that
larger models are not necessarily more truthful.

\paragraph{Factuality and calibration.}
\citet{kadavath2022language} study whether models know what they
know, finding that well-calibrated models can identify their own
uncertainty. \citet{mielke2022reducing} examine confidence
calibration in language model generations.

\paragraph{Mechanistic interpretability.}
\citet{elhage2021mathematical} introduce a mathematical framework
for understanding transformer circuits. \citet{olsson2022context}
identify induction heads as a key mechanism in in-context learning.
Our work applies interpretability tools to the specific problem
of hallucination, examining which internal components fail
during factual recall.

\paragraph{Knowledge in transformers.}
\citet{geva2021transformer} demonstrate that feed-forward layers
in transformers function as key-value memories storing factual
knowledge. This finding motivates our examination of both
attention and feed-forward activations during hallucination.

% ── 3. Method ─────────────────────────────────────────────────────
\section{Method}

\subsection{Models}

We conduct experiments on two models:

\textbf{Custom transformer (806K parameters).} A GPT-style
decoder-only transformer with 4 blocks, 4 attention heads, and
$d_{\text{model}} = 128$. Trained from scratch on a curated
factual corpus of capital cities and geographic facts.
This model serves as a controlled environment where training
data is fully known.

\textbf{GPT-2 (124M parameters).} The standard GPT-2 model
\citep{radford2019language} pretrained on 40GB of internet text.
This model serves as a validation environment representing
production-scale language models.

\subsection{Evaluation Framework}

For each prompt, we record:
\begin{enumerate}
    \item The top-10 predicted tokens and their probabilities
    \item Attention weights in the final transformer block
    averaged across all heads
    \item Hidden state projections at each transformer block,
    obtained by passing the residual stream through the model's
    final layer norm and unembedding matrix
    \item The layer at which the correct answer probability peaks
    \item The layer at which suppression first occurs
\end{enumerate}

\subsection{Hallucination Taxonomy}

We classify each incorrect prediction into one of three types:

\textbf{Type 1 --- Relation Dropout.}
The model fails to attend to the semantic relation token
(e.g., ``capital'') in the final block. Relation attention
falls below a threshold of 0.05.
Characteristic of small, undertrained models.

\textbf{Type 2a --- Last-Layer Suppression.}
The correct answer appears in the top-10 predictions but is
not the top-1. The model ``knows'' the answer but a competing
structural pattern overrides it in the final block.

\textbf{Type 2b --- Knowledge Gap.}
The correct answer does not appear in the top-10 predictions.
The model has not encoded the relevant fact during training.

% ── 4. Experiments ────────────────────────────────────────────────
\section{Experiments}

\subsection{Experiment 1: Relation Dropout in Small Models}

We trained a 806K-parameter transformer on a factual corpus
and tested it on 6 capital city prompts of the form
``The capital of [country] is''.
For each prompt, we measured the average attention weight
directed to the relation token ``capital'' in the final
transformer block across all attention heads.

\subsection{Experiment 2: GPT-2 Validation}

We tested GPT-2 on 35 factual prompts across 5 categories:
capitals, inventors, authors, science facts, and historical
dates. For each prompt we classified the hallucination type
using our taxonomy.

\subsection{Experiment 3: Layer-by-Layer Suppression Analysis}

For 5 prompts exhibiting Type 2a suppression, we projected
the hidden state at each of GPT-2's 12 transformer blocks
through the final layer norm and unembedding matrix to obtain
per-layer vocabulary probabilities. This allowed us to track
the probability of both the correct answer and the competing
token at each layer.

\subsection{Experiment 4: Large-Scale Validation}

Using an NVIDIA RTX 4060 Laptop GPU, we ran 20,000 factual
prompts through GPT-2 across 7 knowledge categories. Prompts
were generated from 50 country-capital pairs (with 3 phrasing
variants each), 15 scientist-discovery pairs, 10 historical
dates, 10 science facts, and 10 author-work pairs.

% ── 5. Results ────────────────────────────────────────────────────
\section{Results}

\subsection{Relation Dropout (Experiment 1)}

Table \ref{tab:relation_dropout} shows the relation attention
scores and hallucination outcomes for our small model.

\begin{table}[h]
\centering
\caption{Relation Dropout in 806K-parameter model.
All 6 hallucination cases show relation attention below 0.05.}
\label{tab:relation_dropout}
\begin{tabular}{lccc}
\toprule
\textbf{Prompt} & \textbf{Predicted} &
\textbf{Rel. Attn} & \textbf{Result} \\
\midrule
...capital of France is  & the    & 0.037 & Hallucination \\
...capital of Germany is & paris  & 0.028 & Hallucination \\
...capital of Italy is   & in     & 0.033 & Hallucination \\
...capital of Japan is   & the    & 0.033 & Hallucination \\
...capital of China is   & the    & 0.083 & Hallucination \\
...capital of Spain is   & madrid & 0.082 & \textbf{Correct} \\
\bottomrule
\end{tabular}
\end{table}

The single correct prediction maintains relation attention above
0.08. Four of five hallucination cases show relation attention
below 0.05, consistent with our Relation Dropout hypothesis.

\subsection{GPT-2 Validation (Experiment 2)}

Table \ref{tab:gpt2} shows results by category for GPT-2.

\begin{table}[h]
\centering
\caption{GPT-2 hallucination taxonomy across 5 categories
(35 prompts). Zero Type 1 cases observed.}
\label{tab:gpt2}
\begin{tabular}{lcccc}
\toprule
\textbf{Category} & \textbf{Correct} &
\textbf{Type 1} & \textbf{Type 2a} & \textbf{Type 2b} \\
\midrule
Capitals   & 4/15 & 0 & 7  & 4 \\
Inventors  & 0/5  & 0 & 3  & 2 \\
Authors    & 0/5  & 0 & 0  & 5 \\
Science    & 2/5  & 0 & 0  & 3 \\
History    & 1/5  & 0 & 1  & 3 \\
\midrule
\textbf{Total} & \textbf{7/35} & \textbf{0} &
\textbf{11} & \textbf{17} \\
\bottomrule
\end{tabular}
\end{table}

\subsection{Last-Layer Suppression (Experiment 3)}

Layer-by-layer analysis reveals a consistent pattern across
all tested prompts. The probability of the correct answer
remains near zero through blocks 1--9, surges in blocks 10--11,
then collapses in block 12 as the structural pattern token
(typically ``the'') takes over.

\begin{table}[h]
\centering
\caption{Peak factual layer and suppression layer for 4 prompts.
Block 12 suppresses in every case.}
\label{tab:layers}
\begin{tabular}{lccc}
\toprule
\textbf{Prompt} & \textbf{Peak Layer} &
\textbf{Peak Prob} & \textbf{Suppression} \\
\midrule
...capital of France is  & Block 10 & 0.182 & Block 12 \\
...capital of Germany is & Block 11 & 0.347 & Block 12 \\
...capital of Japan is   & Block 11 & 0.461 & Block 12 \\
The Berlin Wall fell in  & Block 10 & 0.128 & Block 12 \\
\midrule
\textbf{Average} & \textbf{10.5} & -- & \textbf{12.0} \\
\bottomrule
\end{tabular}
\end{table}

We note that ``The Berlin Wall fell in 1989'' is predicted
correctly despite Block 12 suppression, because the factual
signal (0.128) exceeds the suppressor token's probability
(0.045). We term this the \emph{survival threshold}.

\subsection{Large-Scale Validation (Experiment 4)}

Table \ref{tab:largescale} shows results across 20,000 prompts.

\begin{table}[h]
\centering
\caption{Large-scale results across 20,000 prompts on GPT-2.}
\label{tab:largescale}
\begin{tabular}{lcc}
\toprule
\textbf{Hallucination Type} & \textbf{Count} & \textbf{\%} \\
\midrule
Correct              &    954 &  4.8\% \\
Type 1 (Dropout)     &  2,946 & 14.7\% \\
Type 2a (Suppression)&  2,481 & 12.4\% \\
Type 2b (Gap)        & 13,619 & 68.1\% \\
\midrule
\textbf{Total}       & \textbf{20,000} & \textbf{100\%} \\
\bottomrule
\end{tabular}
\end{table}

The average peak factual layer across all 20,000 prompts is
\textbf{11.1}, and the average suppression layer is
\textbf{12.0}, confirming Last-Layer Suppression at scale.

We also observe strong \textbf{prompt sensitivity}: alternative
phrasings of the same fact dramatically reduce accuracy.
The prompt ``The capital of France is'' achieves 8.1\% accuracy
on capitals, while ``France's capital city is'' achieves 0.0\%,
despite encoding the same factual query.

% ── 6. Discussion ─────────────────────────────────────────────────
\section{Discussion}

\paragraph{Relation Dropout vs Last-Layer Suppression.}
The two named phenomena operate at different scales and through
different mechanisms. Relation Dropout reflects a failure of
attention routing: the model fails to identify which words are
semantically relevant. Last-Layer Suppression reflects a failure
of output selection: the model identifies the correct answer
internally but cannot output it against a stronger structural
pattern.

\paragraph{The survival threshold.}
Our results suggest that facts with high training frequency
can survive Last-Layer Suppression. ``The Berlin Wall fell in
1989'' and ``Water is made of hydrogen and oxygen'' both survive
because their factual signal exceeds the suppressor token.
This implies that hallucination rate is partly a function of
training data frequency, independent of model size.

\paragraph{Prompt sensitivity.}
The dramatic accuracy drop between phrasing variants of the
same factual query (8.1\% vs 0.0\% for capitals) suggests that
GPT-2's factual knowledge is highly form-dependent. The model
has not learned facts in an abstract sense but has learned
specific surface-form patterns.

\paragraph{Implications for detection.}
HallScan operationalises our findings as a practical tool.
By examining relation attention in the final block and
tracking the peak factual layer, HallScan can classify
hallucination risk without access to the correct answer.

% ── 7. Limitations ────────────────────────────────────────────────
\section{Limitations}

Our study has several limitations that future work should address.

First, our small model is trained on a deliberately constrained
corpus, which may not reflect the diversity of real-world
training data. The Relation Dropout threshold of 0.05 was
determined empirically and may not generalise across
architectures.

Second, our large-scale experiment repeats a fixed set of
210 base prompts to reach 20,000 examples. While this tests
consistency of findings, it does not provide 20,000 independent
factual queries. Future work should construct larger independent
prompt sets.

Third, we test only GPT-2 at the production-model scale.
Validation on larger models (GPT-3, LLaMA, Claude) is needed
to confirm whether Last-Layer Suppression is specific to
GPT-2's architecture or a general phenomenon.

Fourth, our taxonomy is based on observable output behaviour
and attention patterns. Deeper mechanistic analysis using
causal intervention methods \citep{elhage2021mathematical}
would strengthen the causal claims.

% ── 8. Conclusion ─────────────────────────────────────────────────
\section{Conclusion}

We have demonstrated that language model hallucinations follow
consistent, detectable internal patterns. We named two phenomena:
Relation Dropout in small models and Last-Layer Suppression
in GPT-2. We validated these findings across 20,000 prompts
and released HallScan and HallBench to support future research.

We believe internal activation analysis represents a promising
direction for hallucination detection that is complementary
to existing output-level approaches. The patterns we identify
suggest that the mechanisms of hallucination are not random
but systematic --- and therefore potentially correctable.

Code, data, and models are available at:\\
\url{https://github.com/TrazeMaG/hallucination-fingerprints}\\
HallBench dataset: \url{https://huggingface.co/datasets/Trazemag/hallbench}

% ── References ────────────────────────────────────────────────────
\bibliographystyle{plainnat}

\begin{thebibliography}{10}

\bibitem[Brown et al.(2020)]{brown2020language}
Brown, T., Mann, B., Ryder, N., et al. (2020).
Language models are few-shot learners.
\emph{Advances in Neural Information Processing Systems}, 33.

\bibitem[Elhage et al.(2021)]{elhage2021mathematical}
Elhage, N., Nanda, N., Olsson, C., et al. (2021).
A mathematical framework for transformer circuits.
\emph{Transformer Circuits Thread}.

\bibitem[Geva et al.(2021)]{geva2021transformer}
Geva, M., Schuster, R., Berant, J., \& Levy, O. (2021).
Transformer feed-forward layers are key-value memories.
\emph{EMNLP 2021}.

\bibitem[Ji et al.(2023)]{ji2023survey}
Ji, Z., Lee, N., Frieske, R., et al. (2023).
Survey of hallucination in natural language generation.
\emph{ACM Computing Surveys}, 55(12).

\bibitem[Kadavath et al.(2022)]{kadavath2022language}
Kadavath, S., Conerly, T., Askell, A., et al. (2022).
Language models (mostly) know what they know.
\emph{arXiv:2207.05221}.

\bibitem[Lin et al.(2022)]{lin2022truthfulqa}
Lin, S., Hilton, J., \& Evans, O. (2022).
TruthfulQA: Measuring how models mimic human falsehoods.
\emph{ACL 2022}.

\bibitem[Mielke et al.(2022)]{mielke2022reducing}
Mielke, S. J., Szlam, A., Dinan, E., \& Boureau, Y. L. (2022).
Reducing conversational agents' overconfidence
through linguistic calibration.
\emph{TACL}, 10.

\bibitem[Olsson et al.(2022)]{olsson2022context}
Olsson, C., Elhage, N., Nanda, N., et al. (2022).
In-context learning and induction heads.
\emph{Transformer Circuits Thread}.

\bibitem[Radford et al.(2019)]{radford2019language}
Radford, A., Wu, J., Child, R., et al. (2019).
Language models are unsupervised multitask learners.
\emph{OpenAI Blog}.

\end{thebibliography}

\end{document}